%----------------------------------------------------------------------------
\chapter{Felhasznált eszközök}
\label{sec:LatexTools}

Munkámat a JVM-ökoszisztéma eszközeivel végeztem, szem előtt tartva a kiírásban kulcsfontosságú elemként rögzített modernitást, a kompatibilitást és a hordozhatóságot. Ennek megfelelően ahol lehetséges, kerültem a jövőbeli eltávolításra kijelölt programelemek használatát, illetve az instabil verzióktól függést.


\section{Kotlin \raisebox{-1ex}{\hspace{1cm}\includegraphics[height=8mm, keepaspectratio]{images/kotlin_logo.png}}}

A Kotlin multiplatform, statikusan típusos, több paradigmát támogató programozási nyelv. Specifikációját a JetBrains rendszeresen frissített, publikus, bárki által megtekinthető formában teszi közzé.~\cite{KotlinSpec} Az első prototípus verzió 2011-ben került kiadásra, és mindössze egy évvel később nyílt forráskódú projektté alakult. Az első éles kiadását 2016-ban érte el a nyelv, 2017-től hivatalosan támogatott az Android mobil operációs rendszerre írt alkalmazások terén, 2019 óta pedig az Android elsődleges nyelve.~\cite{KotlinPast}

Népszerűségét elsősorban az intuitív szintaxisának és a JVM-mel\footnote{JVM: \@Java Virtual Machine, jelenleg az Oracle tulajdonában lévő, szabvánnyal specifikált futtatómotor, elsősorban a Java nyelven írt programok futtatásához } való már-már tökéletes kompatibilitásának köszönheti. Bármilyen Java-ban írt kódrészlet felhasználható külső referenciaként egy Kotlin programban, a programozó közösség által készített számtalan Java könyvtár így kiaknázható marad.  Egy projekten belül nincs korlátozás a két nyelv forrásfájljainak megoszlására; érdemes viszont megjegyezni, hogy a kompatibilitás egyirányú. Java kódból nem lehet visszahívni Kotlin kódba.

A megírt programok fordítási és optimalizálási fázisokon mennek keresztül, mielőtt futtatható állományok vagy eszköztárak készülnek belőlük. A Kotlin fordítómotorjával készített binárisok közel azonos teljesítményűek hagyományos JVM megfelelőikhez képest. A JVM-en kívül más fordítási célok is léteznek:
\begin{itemize}
    \item Kotlin/Native (2023 óta stabil, számos operációs rendszerre biztosít natív, futtató virtuális gépet nem igénylő megoldást)
    \item Kotlin/JS (a megírt kódot JavaScript-re fordítjuk, így egy fordítást igénylő nyelvből interpretált kódot készíthetünk; webböngészőhöz ideális)
\end{itemize}

A nyelv alapvetően egyszerű kódstruktúrájú, jól értelmezhető és gyorsan írható. Emellett a frusztráló, sok esetben hátráltató szintaktikai és programtechnikai elemek közül sokat elhagytak a fejlesztők -- erre és a funkcionális paradigma által nyújtott lehetőségekre a későbbiekben számos példát fogok mutatni kódrészleteken keresztül.

Szakdolgozatom elkészítéséhez a legfrissebb, 2.0.21-es verziót használtam a JetBrains saját, IntelliJ Idea néven futó programján keresztül. A felhasznált JVM változat a legújabb stabilnak tekinthető 23-as főverzió Oracle OpenJDK implementációja; a szakdolgozat során elkészült program által garantált abszolút minimális verzió a 2023 szeptemberében kiadott 21-es LTS\footnote{Long Term Support, a szoftverfejlesztés körében a kiadó által az adott verzióval kapcsolatban vállalt hosszabb idejű, több éves támogatást jelent} változat, ennek a főágú támogatása 2028-ban ér véget.~\cite{JavaRoadmap}

\section{OpenGL \raisebox{-1ex}{\hspace{1cm}\includegraphics[height=8mm, keepaspectratio]{images/opengl_logo.png}}}

Népszerű alacsony szintű grafikus könyvtár, amivel akár pixel szinten címezhető az éppen aktuális futtató felület. Első publikációja 1992-re tehető, a legújabb 4.6-os változat 2017-es dátumával ez a keretrendszer az egyetlen kivétel a modernitási törekvéseim alól.~\cite{OpenglHistory} Robosztussága és széles körű támogatottsága erős fegyvertény, hiszen szinte bármilyen nyelven és operációs rendszeren elérhetőek a kívánt grafikus függvények. Az OpenGL szabványnak megfelelő alkalmazásoknak ablakkezelő rendszertől függetlenül kiszámítható megjelenést kölcsönöz a szoftver.

\section{libGDX \raisebox{-0.1ex}{\hspace{1cm}\includegraphics[height=6mm, keepaspectratio]{images/libgdx_logo.png}}}

A libGDX nevét szerteágazó felhasználási lehetőségeiből kapta: "game development" és "effects" könyvtár Java alapú projektekhez. Segítségével az alacsony szintű OpenGL műveletek a JVM-környezetekre jellemző módon hívhatók meg, így meglepően jól illeszkedhet a grafikus felület kódja az alkalmazás többi részéhez.

Elődje kifejezetten Android alkalmazások megjelenítését segítette elő, 2010 óta közösségi támogatású, nyílt forráskódú eszközként fejlesztik.~\cite{LibgdxHistory} A projekt jelenlegi verziószáma 1.13.0; asztali környezetekben (Windows, macOS, Linux), mobil operációs rendszereken (Android, iOS) és webböngészőkben egyaránt prezentálhatóak az elkészült játékok.

\section{Exposed {\hspace{1cm}\includegraphics[height=8mm, keepaspectratio]{images/exposed_logo.png}}}

Az adatbázisok, azon belül az SQL\footnote{Structured Query Language, relációs adatbázisok kezeléséhez használt szakterület-specifikus nyelv} világa régi alapokon fekszik. Bár az elmúlt évtized is bővelkedik az újításokban - elsősorban fejlesztőbarátabb eszközök létrejöttével - a szabvány közel 40 éves története erősen kihat a mai felhasználásra.

A nyelv egyes dialektusai tulajdonképpen lingvisztikai tájszólásokhoz hasonlóak; az egyik típus "beszélője" a másikat is érteni fogja, de nem biztos, hogy "megszólalni" is helyesen tud. Ezeket az apró eltéréseket számos ORM\footnote{Object Relational Mapping, az adatbázisok és memóriában tárolt objektumok közötti átalakítás technikája} könyvtár kísérelte már meg áthidalni. Talán a legismertebb példák a Hibernate (Java), EntityFramework (.NET/\texttt{C\#}) és az Active Record (Ruby on Rails). A Kotlint fejlesztő JetBrains preferált megoldása a JDBC\footnote{Java Database Connectivity, kifejezetten a Java-hoz írt interfész réteg, célja minél több adatbázis formátum támogatása}-re épülő Exposed, ami tükrözi a Kotlin nagy szabadságot adó viselkedését. Hagyományos objektumorientált módon, avagy a nyers SQL-re nagyban hasonlító szabályokat használó, egyedi nyelvvel érhetjük el a tárolt adatokat.~\cite{ExposedDocs}

Nagy előny a hordozhatóság szempontjából az adatbázissal szembeni függetlenség: a támogatott motorok közötti váltás csupán egy konfigurációs érték módosítását igényli, így a mobilon jellemző SQLite-tól kezdve a szervereken átívelő redundáns Oracle Database implementációkig bármilyen környezetre alkalmas végterméket tudunk kiadni.

% TODO
\section{Ktor {\hspace{1cm}\includegraphics[height=8mm, keepaspectratio]{images/ktor_logo.png}}}